%! Language = Russian
%! TEX program = xelatex

\documentclass[12pt, a4paper]{article}

%  1. Язык и шрифты
\usepackage{fontspec}
\setmainfont{Times New Roman}
\usepackage[russian]{babel}

%  2. Поля и интервалы
\usepackage[left=3cm, right=1.5cm, top=2cm, bottom=2cm]{geometry}
\usepackage{setspace}
\onehalfspacing

%  3. Абзацы
\usepackage{indentfirst}
\setlength{\parindent}{1.25cm}

%  4. Заголовки (пакет titlesec)
\usepackage{titlesec}
\usepackage{titletoc}

% Настройка нумерованных разделов
\titleformat{\section}[block]{\hspace{1.25cm}\bfseries\large}{\thesection}{1em}{}
\titleformat{\subsection}[block]{\hspace{1.25cm}\bfseries\normalsize}{\thesubsection}{1em}{}

% Команда для структурных элементов
\newcommand{\structsection}[1]{
    \newpage
    \addcontentsline{toc}{section}{#1}
    \begin{center}
    \MakeUppercase{#1}
    \end{center}
    \par\vspace{1em}
}

%  5. Оформление сносок
\usepackage[hang]{footmisc}
\renewcommand{\footnotesize}{\fontsize{10pt}{12pt}\selectfont}
\setlength{\footnotemargin}{1.25cm}

%  6. Приложения
\newcommand{\appendixsection}[2]{
    \newpage
    \begin{center}
    ПРИЛОЖЕНИЕ #1 \\
    \vspace{1em}
    \textbf{#2}
    \end{center}
    \addcontentsline{toc}{section}{ПРИЛОЖЕНИЕ #1. #2}
    \par\vspace{1em}
}

\begin{document}

    \begin{titlepage}
        \begin{center}
            Автономная некоммерческая организация высшего образования \\
            «Русский университет метатехнологий» \\
            \vspace{1em}
            Кафедра программных систем
        \end{center}

        \vspace{2em}
        \noindent Направление подготовки: 09.03.04 Программная инженерия \\
        Профиль подготовки: Проектирование и разработка программных систем

        \vspace{5em}
        \begin{center}
            ПОЯСНИТЕЛЬНАЯ ЗАПИСКА К КУРСОВОМУ ПРОЕКТУ \\
            \vspace{1em}
            на тему: \\
            «Название работы»
        \end{center}

        \vspace{5em}
        \begin{flushright}
            \begin{minipage}{0.45\textwidth}
                Выполнил(а): \\
                студент(ка) группы \rule{2cm}{0.15mm} \\
                очной формы обучения \\
                \rule{3cm}{0.15mm} / Фамилия И.О. \\
                {\footnotesize подпись \hspace{1cm} расшифровка} \\

                \vspace{2em}
                Руководитель: \\
                должность, уч. степень \\
                \rule{3cm}{0.15mm} / Фамилия И.О. \\
                {\footnotesize подпись \hspace{1cm} расшифровка}
            \end{minipage}
        \end{flushright}

        \vfill
        \begin{center}
            Йошкар-Ола \\
            \the\year
        \end{center}
    \end{titlepage}

    \setcounter{page}{2}

    \newpage
    \addcontentsline{toc}{section}{СОДЕРЖАНИЕ}
    \begin{center}
        СОДЕРЖАНИЕ
    \end{center}
    \tableofcontents

    \structsection{Введение}
    Текст введения. Здесь описывается актуальность, цель и задачи курсового проекта. Обратите внимание на абзацный отступ, который применяется автоматически ко всем абзацам.

    \newpage


    \section{Аналитическая часть}
    Описание используемых методов при разработке собственного языка программирования\footnote{Пример пояснения, оформленного в виде сноски 10 пт.}.

    \subsection{Обзор существующих решений}
    Текст подраздела. Форматирование подразделов идет с абзацного отступа, полужирным шрифтом без точки в конце.

    \newpage


    \section{Техническая часть}
    Описание особенностей реализации программного продукта. Ссылки на литературу оформляются в квадратных скобках [1].

    \structsection{Заключение}
    В заключении подводятся итоги выполненной работы, описываются возникшие проблемы и степень их проработанности.

    \appendixsection{А}{Листинг кода}
    Здесь могут быть приведены схемы, оригинальные фрагменты кода и т.п.

\end{document}